% SEGMENT OLP-0434-S01
% Part: normal-modal-logic
% Chapter: axioms-systems
% Section: duals

\documentclass[../../../include/open-logic-section]{subfiles}

\begin{document}

\olfileid{nml}{prf}{dua}

\olsection{இருமை \usetoken{P}{formula}}

\begin{defn}\ollabel{def:duals}
  \Ax{T}, \Ax{B}, \Ax{4}, \Ax{5} என்ற !!{formula}s ஒவ்வொன்றுக்கும்,
  கீழொட்டிட்ட வைரத்தால் குறிக்கப்படும் ஓர் \emph{இருமை} பின்வருமாறு உள்ளது:
  \begin{align}
    \tag{\Ax{T_\Diamond}} p & \lif \Diamond p\\
    \tag{\Ax{B_\Diamond}} \Diamond\Box p & \lif p\\
    \tag{\Ax{4_\Diamond}} \Diamond\Diamond p & \lif \Diamond p\\
    \tag{\Ax{5_\Diamond}} \Diamond\Box p & \lif \Box p
    \end{align}
\end{defn}

மேலுள்ள இருமை !!{formula}s ஒவ்வொன்றும், அதற்குரிய !!a{formula} இல்
$p$ க்குப் பதிலாக $\lnot p$ ஐப் பதிலிட்டு, முரண்மாற்றி,
$\lnot\Box\lnot$ ஐ~$\Diamond$ ஆலும் $\lnot\Diamond\lnot$ ஐ~$\Box$ ஆலும்
பதிலிடுவதால் பெறப்படுகிறது. \Ax{D}, அதாவது $\Box!A \lif \Diamond!A$,
இந்தப் பொருளில் தனக்குத்தானே இருமையாகும்.

\begin{prop}\ollabel{prop:dualsys}
  \olref[nml][prf][dua]{def:duals} இலுள்ள ஒவ்வொரு !!a{formula}~$!A$ க்கும்,
  $\Log{K}!A = \Log{K}!A_{\Diamond}$.
\end{prop}
\begin{proof}
  பயிற்சி.
\end{proof}
\begin{prob}
  \olref[nml][prf][dua]{prop:dualsys} ஐ நிறுவுக.
\end{prob}

\end{document}
